\usepackage[T1]{fontenc}
\usepackage[utf8]{inputenc}
\usepackage[sfdefault]{FiraSans}
\usepackage{newtxsf}% sans serif math, to match the text
\usepackage{FiraMono}
\usepackage{tikz}
\usepackage{pgfplots}
\pgfplotsset{compat=1.18}
\usepackage{fontawesome}
\usepackage{csquotes}
\usepackage{morewrites}
\usepackage{upquote}
\usepackage{multicol}
\usepackage{booktabs,array,tabularx}
\usepackage{siunitx}
% all numbers with units go through \qty, \num, and \unit. They follow the surrounding font (bold, italic, text or math),
% thin spaces group digits from 1000 on and separate a number from its unit, an explicit plus sign stays
\sisetup{
   mode=match,
   propagate-math-font=true,
   reset-math-version=false,
   reset-text-family=false,
   reset-text-series=false,
   reset-text-shape=false,
   text-family-to-math=true,
   text-series-to-math=true,
   group-separator={\,},
   group-minimum-digits=4,
   output-decimal-marker={.},
   retain-explicit-plus=true,
   quantity-product={\,},
   per-mode=symbol,
   range-units=single,
   range-phrase={ to },
}
\DeclareSIUnit\euro{\text{\texteuro}}
\usepackage{fancyqr/fancyqr}
\usetikzlibrary{arrows.meta,decorations.pathmorphing,decorations.pathreplacing,fit,backgrounds,matrix,shadows,positioning,calc,shapes.geometric,shapes.misc,fadings,trees,shapes.callouts}
\usepgfplotslibrary{fillbetween,groupplots}
\pgfdeclarelayer{foreground}
\usepackage[glows]{tikzpingus}
\pgfsetlayers{very-background,background,main,middle,foreground}
\def\sbseries{\fontseries{sb}\selectfont}
\def\textsb#1{{\sbseries#1}}

\usepackage{slide-template-uulm/fancybeamer}
\usepackage{slide-template-uulm/fancyuulm}
\setpaths{{slide-template-uulm/}{slide-template-uulm/logos/}{slide-template-uulm/empty-slides/}{xlistings/}{magic-r/}{figures/out/}}

\usepackage[auto]{contour}
\PassOptionsToPackage{normalem}{ulem}
\usepackage{ulem}
\makeatletter
\DeclareRobustCommand*\fancyul[2][\fancyulbackground]{%
   \def\fancyul@background{#1}%
   \def\ULdepth{1pt}%
   \def\ULthickness{.1ex}%
   \contourlength{.35pt}%
   {\color{\fancyulcolor}\uline{\phantom{#2}}}\llap{\contour\fancyul@background{#2}}%
}
\def\fancyulbackground{white}
\colorlet{ul@line@color}{blue!35}
\def\fancyulcolor{ul@line@color}
\def\update@ul{\setbox\z@\hbox{{(j}}\let\ULdepth\p@\edef\ULthickness{\the\dimexpr.235\dp\z@}}
\DeclareRobustCommand*\link{\hyper@normalise\link@do}
\def\link@do#1#2{\leavevmode\hskip-2\p@\href{#1}{\update@ul\hskip2\p@\fancyul{#2}}}
\makeatother
\hypersetup{colorlinks=false,pdfborder=0 0 0}
\usepackage[tikz]{ocgx2}
\urlstyle{same}

\usepackage{code-animation/code-animation}
\usepackage[fakeminted,print]{xlistings/xlistings}
\xlstsetmintedstyle{plain}
\LoadLanguages{yaml}
\usepackage[language=R,eval,graphs,playground,underline=false]{magic-r/magic-r}
% magic-r's play arrow, left of the language badge, both above the first line and out of the way of the code
\makeatletter\def\xlstlangbadgestyle#1{\raisebox{1.3em}{\normalfont\color{xlst@colors@lst@linenumbers}\textsf{\footnotesize#1}}}\makeatother
\mrPlayStyle{filled}
\def\mrPlayOffset{-2.6mm}
\renewcommand\mrPlaceButton[1]{\rlap{\hspace*{\linewidth}\hspace{\mrPlayOffset}\llap{\smash{\raisebox{1.3em}{#1}}}}}
\mode<handout>{\mrSessionSetup{animate=false}}
% a session has no badge and often text right above it: magic-r's button and an 'R' at the right of its
% first line, linked to the session's inputs in flowR's playground (\mrSessionPlayUrl, empty without one)
% a small gap before each input of a session: the prompt makes its line taller, and the fixed depth of every line keeps the gaps equal
\let\LecturePlainPrompt\mrPrompt
\renewcommand\mrPrompt[1]{\ifstrequal{#1}{>}{\rule{0pt}{\dimexpr.7\baselineskip+1pt\relax}}{}\LecturePlainPrompt{#1}}
\makeatletter
\lst@AddToHook{EveryPar}{\rule[-.3\baselineskip]{0pt}{0pt}}
\lst@AddToHook{EveryPar}{\ifx\mrSessionPlayUrl\@empty\else
   % an hbox within \smash: the link takes its height from the enclosing box, which \smash would make zero
   \rlap{\hspace*{\linewidth}\llap{\smash{\hbox{\expandafter\href\expandafter{\mrSessionPlayUrl}{\mrPlayButton\hspace{1mm}%
      \normalfont\color{xlst@colors@lst@linenumbers}\textsf{\footnotesize R}}}}}}%
   \global\let\mrSessionPlayUrl\@empty\fi}
\makeatother
% keywords dark red, packages sand, functions plum, numbers orange, strings green: each kind tells apart at a glance
\definecolor{codeKeyword}{HTML}{7A2430}
\definecolor{codeFunction}{HTML}{5E3A68}
\definecolor{codePackage}{HTML}{8A6A2E}
\definecolor{codeNumber}{HTML}{A85A12}
\definecolor{codeString}{HTML}{2E7D32}
\lstcolorlet{keywordA}{codeKeyword}
\lstcolorlet{keywordB}{codePackage}
\lstcolorlet{keywordC}{codeFunction}
\lstcolorlet{numbers}{codeNumber}
\lstcolorlet{literals}{codeString}
% named arguments ('n' in 'head(x, n = 3)') only slightly colored
\definecolor{codeArgument}{HTML}{7A6A80}
\mrSetStyle{argument}{\color{codeArgument}}
% a little air between the lines of code, the same between every two lines (no \lineskip where a line is tall)
\makeatletter\lst@AddToHook{Init}{\baselineskip=1.18\baselineskip\lineskiplimit=-\maxdimen}\makeatother
\xlstmintedwithlangbadge
\makeatletter
\def\ca@Strut{\vphantom{\ttfamily\strut}}
\makeatother

\usepackage[style=numeric,sorting=nyt,defernumbers=true,backend=biber,maxbibnames=4,maxcitenames=2,sortcites=true]{biblatex}
\addbibresource{references.bib}
% a DOI is enough, the URL would only repeat it
\AtEveryBibitem{\iffieldundef{doi}{}{\clearfield{url}}}
\AtEveryBibitem{\clearfield{urlyear}\clearfield{urlmonth}\clearfield{urlday}\clearfield{isbn}\clearfield{issn}\clearfield{note}\clearfield{version}\clearfield{eprint}\clearfield{pages}\clearlist{location}}
% the three kinds of references (keywords in references.bib): a colour for their citations, a marker in the bibliography
\newcommand\CiteKind[3]{\csdef{cite@#1@colour}{#2}\csdef{cite@#1@icon}{#3}}
\CiteKind{fundamentals}{plotblue}{\faStar}
\CiteKind{resources}{Gentoo}{\faSearch}
\CiteKind{niceToKnow}{Adelie}{\faLightbulbO}
\newcommand\IfCiteKind[1]{%
   \ifkeyword{fundamentals}{#1{fundamentals}}{\ifkeyword{resources}{#1{resources}}{#1{niceToKnow}}}}
\newcommand\CiteColoured[2]{\textcolor{\csuse{cite@#1@colour}!85!black}{#2}}
\DeclareFieldFormat{labelnumber}{\IfCiteKind{\CiteMark{#1}}}
\def\CiteMark#1#2{\CiteColoured{#2}{#1}}

% every citation of a frame is listed at its bottom with author, year, and venue
\newif\ifListingCites
\AtEveryCitekey{\ifListingCites\else\xifinlist{\thefield{entrykey}}\FrameCites{}{\listxadd\FrameCites{\thefield{entrykey}}}\fi}
\newbibmacro{venue}{%
   \iffieldundef{usera}{\iffieldundef{journaltitle}{\iffieldundef{booktitle}{\iflistundef{publisher}{\iflistundef{organization}%
      {\iffieldundef{eprinttype}{}{\printfield{eprinttype}}}{\printlist{organization}}}{\printlist{publisher}}}%
      {\printfield{booktitle}}}{\printfield{journaltitle}}}{\printfield{usera}}}
% \FrameCite{key}: lists key at the bottom of the frame without a mark in the text
\DeclareCiteCommand{\FrameCite}{}{}{}{}
% a reference at the bottom of a frame links to its DOI, else to its URL, else to its entry in the bibliography
\newbibmacro{footlink}[1]{%
   \iffieldundef{doi}{\iffieldundef{url}{\bibhyperref{#1}}{\href{\thefield{url}}{#1}}}{\href{https://doi.org/\thefield{doi}}{#1}}}
\DeclareCiteCommand{\FootCite}{}{%
   \usebibmacro{footlink}{\mbox{\mkbibbrackets{\printfield{labelnumber}}\addnbspace
   \printnames{labelname}\setunit{\addspace}\iffieldundef{year}{\printfield{urlyear}}{\printfield{year}}%
   \setunit{\addcomma\addspace}\textit{\usebibmacro{venue}}}}}{\hskip 1em\relax}{}
\DeclareCiteCommand{\FootCiteLong}{}{%
   \usebibmacro{footlink}{\mbox{\mkbibbrackets{\printfield{labelnumber}}\addnbspace
   \printnames{labelname}\setunit{\addspace}\iffieldundef{year}{\printfield{urlyear}}{\printfield{year}}%
   \setunit{\addcolon\addspace}\printfield[titlecase]{title}\setunit{\addcomma\addspace}\textit{\usebibmacro{venue}}}}}{\hskip 1em\relax}{}
% a single reference on a frame has room for its title
% (the short form, when the long one would run into the R code links or off the page)
\newdimen\FootCiteRoom
\newcommand\FootCiteLongItem[1]{\CiteSep
   \sbox0{\FootCiteLong{#1}}\ifdim\wd0>\FootCiteRoom\FootCite{#1}\else\FootCiteLong{#1}\fi
   \def\CiteSep{\penalty0\hskip 1em\relax}}
\newcommand\FootCiteItem[1]{\CiteSep\FootCite{#1}\def\CiteSep{\penalty0\hskip 1em\relax}}
% a superscript cite does not push its line away from the one above
% an hbox within \smash: the link takes its height from the enclosing box, which \smash would make zero
\def\supercite#1{\textsuperscript{\smash{\hbox{\color{gray}\cite{#1}}}}}
% [1,2] instead of [1, 2]
\renewcommand*\multicitedelim{\addcomma}
% one space after a full stop, as everywhere else
\frenchspacing

% \NoteRef{key}: speaker notes that live in notes/notes.tex (private), if that file exists
% Speaker notes kept outside the sources: notes/notes.tex holds one DefineNote entry per note (key and text), the frames hold a NoteRef with the key.
% Without notes/notes.tex (the public copy) every \NoteRef and \NextNoteRef expands to nothing.
% tools/extract-notes moves the notes out of the sources.
\makeatletter
\newcommand\DefineNote[2]{\long\expandafter\gdef\csname lecture@note@#1\endcsname{#2}}
% \NoteRef{key} is \note{text}, if the key is defined
\newcommand\NoteRef[1]{\ifcsname lecture@note@#1\endcsname
   \expandafter\expandafter\expandafter\note\expandafter\expandafter\expandafter{\csname lecture@note@#1\endcsname}%
\else\LectureNoteMissing{#1}\fi}
% \NextNoteRef{key} is \NextNote{text}, the note of the next \Quiz or \Takeaways frame, if the key is defined
\newcommand\NextNoteRef[1]{\ifcsname lecture@note@#1\endcsname
   \expandafter\expandafter\expandafter\NextNote\expandafter\expandafter\expandafter{\csname lecture@note@#1\endcsname}%
\else\LectureNoteMissing{#1}\fi}
% a key without text only warns, and only if notes/notes.tex is there (the public copy has none)
\newif\iflecture@notesfile
\newcommand\LectureNoteMissing[1]{\iflecture@notesfile\PackageWarningNoLine{lecture}{No text for the note #1 in notes/notes.tex}\fi}
% the notes are read when the document starts, so that their text is read with the same catcodes as in the frames
\AtBeginDocument{\IfFileExists{notes/notes.tex}{\lecture@notesfiletrue\input{notes/notes.tex}}{}\IfFileExists{notes/title.tex}{\input{notes/title.tex}}{}}
\makeatother

% speaker notes: pdfpc annotations in the animated and the notes build, beamer's own note pages only in the
% notes build, and neither in the handouts
\usepackage{letltxmacro}% the disabled pdfpc package uses \LetLtxMacro before it loads letltxmacro
\ifdefined\LectureAnimated\else\ifdefined\LectureNotes\else
   \PassOptionsToPackage{disable}{beamer-latex-pdfpc-notes}
\fi\fi
\usepackage{beamer-latex-pdfpc-notes}
% the plain text of a note for pdfpc: its markers become lines and labels
\ExplSyntaxOn
\tl_new:N \l__lecture_pdfpc_tl
\LetLtxMacro\LecturePdfpcPlain\note
\NewDocumentCommand\LectureNoteForPdfpc{m}
   {
      \tl_set:Nn \l__lecture_pdfpc_tl {#1}
      \regex_replace_all:nnN { \c{\#} } { \cO\# } \l__lecture_pdfpc_tl
      % pdfpc renders markdown: the main message in bold, then the points as a list, then one paragraph per block
      \regex_replace_once:nnN { \A MAIN : \s (.*?) (\s (?: POINT | SUB | NEXT | TIP | ASK | SHOW | EXTRA | DETOUR | ANECDOTE | READ ) : ) }
         { **\1** \c{\\} \2 } \l__lecture_pdfpc_tl
      \regex_replace_all:nnN { \s POINT : \s } { \c{\\}-\  } \l__lecture_pdfpc_tl
      \regex_replace_all:nnN { \s SUB : \s } { \c{\\}\ \ \ \ -\  } \l__lecture_pdfpc_tl
      \regex_replace_all:nnN { \s ASK : \s } { \c{\\}\c{\\}**Ask:**\  } \l__lecture_pdfpc_tl
      \regex_replace_all:nnN { \s SHOW : \s } { \c{\\}\c{\\}**Show:**\  } \l__lecture_pdfpc_tl
      \regex_replace_all:nnN { \s TIP : \s } { \c{\\}\c{\\}Tip:\  } \l__lecture_pdfpc_tl
      \regex_replace_all:nnN { \s EXTRA : \s } { \c{\\}\c{\\}Extra:\  } \l__lecture_pdfpc_tl
      \regex_replace_all:nnN { \s DETOUR : \s } { \c{\\}\c{\\}**Detour:**\  } \l__lecture_pdfpc_tl
      \regex_replace_all:nnN { \s ANECDOTE : \s } { \c{\\}\c{\\}Anecdote:\  } \l__lecture_pdfpc_tl
      \regex_replace_all:nnN { \s NEXT : \s } { \c{\\}\c{\\}Next:\  } \l__lecture_pdfpc_tl
      \regex_replace_all:nnN { \s READ : \s } { \c{\\}\c{\\}Read:\  } \l__lecture_pdfpc_tl
      \regex_replace_all:nnN { \s ANSWER : \s } { \c{\\}\c{\\}Answer:\  } \l__lecture_pdfpc_tl
      \regex_replace_all:nnN { \s SOURCE : \s } { \c{\\}Source:\  } \l__lecture_pdfpc_tl
      \exp_args:NV \LecturePdfpcPlain \l__lecture_pdfpc_tl
   }
\ExplSyntaxOff
\ifdefined\LectureAnimated\ifdefined\LectureNotes\else\let\note\LectureNoteForPdfpc\fi\fi
\ifdefined\LectureNotes
   \input{setup/notes}
\else
   \newcommand\NoteTimetable{}
   \newcommand\NoteBoard[1]{}
\fi
\setcounter{tocdepth}{1}
